We propose a unified objective function, termed CATMIL, that augments the base segmentation loss with two auxiliary supervision terms operating at different levels. The first term, Component-Adaptive Tversky, reweights voxel contributions based on connected components to balance the influence of lesions of different sizes. The second term, based on Multiple Instance Learning, introduces lesion-level supervision by encouraging the detection of each lesion instance. These terms are combined with the standard nnU-Net loss to jointly optimize voxel-level segmentation accuracy and lesion-level detection. We evaluate the proposed objective on the MSLesSeg dataset using a consistent nnU-Net framework and 5-fold cross-validation. The results show that CATMIL achieves the most balanced performance across segmentation accuracy, lesion detection, and error control. It improves Dice score (0.7834) and reduces boundary error compared to standard losses. More importantly, it substantially increases small lesion recall and reduces false negatives, while maintaining the lowest false positive volume among compared methods. These findings demonstrate that integrating component-level and lesion-level supervision within a unified objective provides an effective and practical approach for improving small lesion segmentation in highly imbalanced settings. All code and pretrained models are available at https://github.com/luumsk/SmallLesionMRI.
